problem:  
Let \((X, d_X, \mu_X)\) be a compact metric measure space satisfying the \(\mathrm{RCD}(K,N)\) condition for some \(K \in \mathbb{R}\), \(1 < N < \infty\). Assume that \(X\) is homeomorphic to a smooth compact manifold \(M\) but has finitely many isolated conical singularities \(p_1, \dots, p_m\), each locally modeled on a metric cone over a compact \(\mathrm{RCD}(N-1, N-1)\) space.  

Let \((Y, d_Y)\) be a compact \(\mathrm{CAT}(0)\) space. For each continuous \(f_0: X \to Y\) that is Lipschitz on the regular part of \(X\), the Korevaar–Schoen energy functional  
\[
E(f) = \int_X |\nabla f|^2 \, d\mu_X
\]
is well-defined.  

**Question:**  
Assume that energy-minimizing harmonic maps \(f: X \to Y\) exist in the Korevaar–Schoen sense. Can such a minimizer develop singular behavior precisely at the conical singularities \(p_i\), or can one prove that \(f\) is necessarily Hölder continuous across each \(p_i\)? In particular, is there a curvature threshold \(K_0 = K_0(N)\) such that for all \(K > K_0\), every harmonic map from such an \(X\) into a compact \(\mathrm{CAT}(0)\) target extends Hölder continuously across the singularities?

Why is it a "good" problem:  
This problem carves out a sharply focused and genuinely open question at the interface of geometric analysis and metric geometry. It no longer concerns mere existence of harmonic maps (largely settled by Korevaar–Schoen and successors) but instead targets their *fine regularity near singular points* in the emerging landscape of spaces with synthetic curvature bounds. Specifically:  

1. **Pinpointing analytic breakdown at singularities:**  
   For smooth manifolds, elliptic regularity yields smooth harmonic maps. In contrast, for \(\mathrm{RCD}\) spaces with conical metrics, the failure of local Euclidean structure makes it unknown whether harmonic minimizers avoid energy concentration at singular points. A result confirming or refuting Hölder regularity would illuminate new analytic phenomena peculiar to synthetic curvature geometries.  

2. **Bridging synthetic geometry and classical PDE theory:**  
   The problem calls for importing tools from PDE regularity theory (monotonicity formulas, blow-up analysis, potential theory) into the framework of \(\mathrm{RCD}\) metric analysis, thus demanding a genuinely cross-disciplinary synthesis.  

3. **Simple statement, deep challenge:**  
   The question—“Do harmonic maps remain regular across conical singularities under lower Ricci curvature bounds?”—is conceptually elementary yet opens a deep analytic frontier bridging smooth and singular curvature theories.  

4. **Potential to influence multiple fields:**  
   A breakthrough would impact the understanding of elliptic regularity on singular spaces, the heat-flow approach to harmonic maps on \(\mathrm{RCD}\) spaces, and applications in metric geometric topology and geometric group theory.

Therefore, this refined problem isolates the true difficulty—the interplay between singular structure and harmonic map regularity—making it both approachable in formulation and profound in mathematical significance.